\chapter{Market Equilibrium and Imperfect Competition}

\section{Competitive Market Equilibrium}

\subsection{Short-Run Equilibrium}

In a perfectly competitive market, buyers and sellers are sufficiently large in number so that no single buyer or seller can affect the market price.

A buyer's demand for a good is the outcome of utility-maximizing behavior, while a seller's supply of that good is the outcome of profit-maximizing behavior.

The demand side of a market is made up of all individual buyers of the good, and the supply side is made up of all sellers of the same good.

Market demand and market supply together determine the market price and the total quantity traded.

There are \(I\) buyers. Let
\[
q_i(p,p^0,m_i)
\]
be individual \(i\)'s demand for good \(q\), where \(p\) is the own price of \(q\), and \(p^0\) denotes the prices of all other goods.

Market demand for \(q\) is
\[
q^d
=
\sum_{i=1}^I q_i(p,p^0,m_i).
\]

Let there be \(J\) sellers.

In the short run, we assume that \(J\) is fixed. Some inputs may also be fixed inputs.

Firm \(j\), one of the existing firms, has a supply function for good \(q\):
\[
q_j^s(p,w).
\]

The short-run market supply function is
\[
q^s(p,w)
=
\sum_{j=1}^J q_j^s(p,w).
\]

A competitive market is in short-run partial equilibrium at price \(p\) when
\[
q^s(p,w)
=
q^d(p).
\]

\paragraph{Exercise 1.}
There are
\[
J=50
\]
identical firms.

A typical firm's production function is
\[
q=x^{0.5}\bar z^{0.5},
\]
where \(\bar z\) is fixed at
\[
\bar z=1.
\]

Input prices are
\[
w_x=4,
\qquad
w_z=1.
\]

Market demand is given by
\[
q^d=\frac{300}{p}.
\]

Derive the short-run market price, market quantity, output per firm, and firm profit.



\subsection{Long-Run Equilibrium}

In the long run, firms are free to leave or enter the industry.

All inputs are variable.

In a long-run equilibrium, there must be no incentive to enter or exit the industry. Thus, we require an additional equilibrium condition: zero profit.

The two equilibrium conditions are
\[
q^s(p)=q^d(p),
\]
and
\[
\pi_j(p,w)=0,
\qquad
j=1,\dots,J.
\]

In the long run, the number of firms is not fixed. Hence, the long-run equilibrium price \(p\) and the long-run number of firms \(J\) must be determined jointly.

\paragraph{Exercise 2.}
The production function is
\[
q=x^{0.5}z^{0.5},
\]
where both \(x\) and \(z\) are variable now.

Again,
\[
w_x=4,
\qquad
w_z=1,
\]
but entry and exit are free of barriers.

Market demand is given by
\[
q^d=\frac{300}{p}.
\]

Derive the long-run market price, market quantity, output per firm, and the number of firms in the industry.

Discuss whether the number of firms \(J^*\) is uniquely determined.

\subsection{Efficiency of the Competitive Equilibrium}

Consumer surplus is a money measure of the gains going to consumers as a result of purchasing a good.

Likewise, producer surplus is the gain going to producers as a result of selling the good.

Producer surplus is the firm's revenue minus its variable cost:
\[
PS
=
p(q)q-tvc(q),
\]
where \(p(q)\) is the inverse market demand function.

In the case of a single consumer and a single producer, the total surplus, defined as the sum of consumer surplus and producer surplus, must be maximized to obtain an efficient outcome.



\[
TS
=
CS+PS
\]

\[
=
\left[
\int_0^q p(x)\,dx-p(q)q
\right]
+
\left[
p(q)q-tvc(q)
\right]
\]

\[
=
\int_0^q p(x)\,dx-tvc(q)
\]

\[
=
\int_0^q
\left[
p(x)-mc(x)
\right]dx.
\]

The first-order condition is
\[
p(q)=mc(q),
\]
where
\[
mc(q)
\]
is marginal cost.

This first-order condition occurs precisely at the competitive market equilibrium when demand is downward-sloping and marginal cost is increasing.

\section{Imperfect Competition}

\subsection{Pure Monopoly}

In pure monopoly, there is a single seller of a product for which there are no close substitutes in consumption, and entry into the market is completely blocked.

Profit is
\[
\pi(q)
=
p(q)q-c(q)
=
r(q)-c(q),
\]
where
\[
r(q)=p(q)q
\]
is total revenue.

Profit is maximized at
\[
q^*
\]
such that
\[
mr(q^*)=mc(q^*).
\]

The equilibrium price is
\[
p^*=p(q^*).
\]

Marginal revenue is
\[
mr(q)
=
\frac{dr(q)}{dq}
=
p(q)+q\frac{dp(q)}{dq}.
\]

Using the elasticity of market demand,
\[
\varepsilon(q)
=
\frac{dq}{dp}\frac{p(q)}{q},
\]
we can rewrite marginal revenue as
\[
mr(q)
=
p(q)
\left[
1+\frac{1}{\varepsilon(q)}
\right].
\]

Since market demand is downward-sloping,
\[
\varepsilon(q)<0,
\]
which implies
\[
mr(q)<p(q).
\]

Because
\[
p'(q)<0,
\]
the monopoly quantity
\[
q^*
\]
is smaller than the competitive equilibrium quantity \(\bar q\), which satisfies
\[
p(\bar q)=mc(\bar q).
\]

From the monopoly first-order condition,
\[
p(q^*)
\left[
1+\frac{1}{\varepsilon(q^*)}
\right]
=
mc(q^*),
\]
we obtain
\[
\frac{
p(q^*)-mc(q^*)
}{
p(q^*)
}
=
-\frac{1}{\varepsilon(q^*)}.
\]

This expression is called the monopolist's markup, or the Lerner index of monopoly power.

Monopoly power increases as market demand becomes more inelastic.



\subsection{Oligopoly: Cournot Competition}

There are \(J\ge 2\) identical firms, and entry by additional firms is blocked.

Each firm has identical cost function
\[
c(q_j).
\]

Market price depends on the total output sold by all firms in the market:
\[
p(q)=p(q_j+q_{-j}),
\]
where
\[
q_{-j}
=
\sum_{k\ne j}q_k.
\]

Firm \(j\)'s profit is
\[
\pi_j(q_1,\dots,q_J)
=
p(q_j+q_{-j})q_j-c(q_j).
\]

The first-order condition is
\[
p(q)
+
p'(q)
\left(
1+\frac{\partial q_{-j}}{\partial q_j}
\right)q_j
-
mc(q_j)
=
0.
\]

In the Cournot--Nash equilibrium, firm \(j\) assumes that
\[
\frac{\partial q_{-j}}{\partial q_j}=0.
\]

Therefore, the first-order condition becomes
\[
p(q)+p'(q)q_j-mc(q_j)=0.
\]

Equivalently,
\[
mr(q_j)-mc(q_j)=0.
\]

If firms are identical, then
\[
q_j=\frac{q}{J}.
\]

Using the elasticity of market demand,
\[
\varepsilon(q)
=
\frac{dq}{dp}\frac{p}{q},
\]
we obtain the markup formula
\[
\frac{p-mc}{p}
=
\frac{1}{J|\varepsilon(q)|}.
\]

This markup is smaller than that of a pure monopolist, but remains positive.

\paragraph{Exercise 3.}
Assume
\[
J=2,
\qquad
p=a-(q_1+q_2),
\]
and
\[
c(q_j)=cq_j.
\]

Derive the equilibrium price, quantity, and profit in competitive, pure monopoly, and Cournot markets, respectively.

\subsection{Oligopoly: Bertrand Competition}

In the Bertrand model, firms compete in their choice of price rather than quantity.

The demand for firm \(j\)'s product depends on its own price and the prices of all other firms:
\[
q_j=q_j(p)=q_j(p_1,\dots,p_J),
\]
where
\[
\frac{\partial q_j}{\partial p_j}<0
\]
and
\[
\frac{\partial q_j}{\partial p_k}>0,
\qquad
k\ne j.
\]

In addition, we assume there is some price
\[
\bar p_j>0
\]
at which demand for firm \(j\)'s product is zero, regardless of other firms' prices.

The profit function is
\[
\pi_j(p)
=
p_jq_j(p)-c(q_j(p)).
\]

The first-order condition is
\[
q_j(p)
+
p_j\frac{\partial q_j(p)}{\partial p_j}
-
mc(q_j(p))
\frac{\partial q_j(p)}{\partial p_j}
=
0.
\]

Equivalently,
\[
q_j(p)
+
\left[
p_j-mc(q_j(p))
\right]
\frac{\partial q_j(p)}{\partial p_j}
=
0.
\]

This can be written as
\[
\frac{\partial q_j(p)}{\partial p_j}
\left[
mr(q_j(p))-mc(q_j(p))
\right]
=
0.
\]

Since
\[
p_j<\bar p_j,
\]
we have
\[
\frac{\partial q_j(p)}{\partial p_j}\ne 0.
\]

Thus,
\[
mr(q_j(p))=mc(q_j(p)).
\]

Therefore,
\[
\frac{
p_j-mc(q_j(p))
}{
p_j
}
=
\frac{1}{|\varepsilon_j(q(p))|},
\]
where
\[
\varepsilon_j(q(p))
\]
is the elasticity of firm \(j\)'s demand.

Assume that
\[
mc(q_j(p))=c.
\]

If firms produce homogeneous outputs, then the goods are perfect substitutes. That is,
\[
|\varepsilon_j(q(p))|=\infty.
\]

The markup formula becomes
\[
\frac{p_j-c}{p_j}
=
\frac{1}{|\varepsilon_j(q(p))|}
=
0.
\]

Hence,
\[
p_j=c.
\]

The markup disappears, and the competitive equilibrium price is maintained.

Because the firm with the lowest price serves the entire market demand at the price it declares, every firm has an incentive to set its price at marginal cost \(c\). As a result, each firm shares market demand equally.

This must be true even when
\[
J=2.
\]

That is, two firms are sufficient to create a competitive market.

In reality, firm products may not be identical, leading to price differentials. Even in the case of homogeneous goods, capacity limitations may prevent the lowest-priced firm from capturing the entire market.

\paragraph{Exercise 4.}
Assume
\[
J=2,
\]
\[
q_j(p_j,p_k)=a-p_j+bp_k,
\qquad
b>0,
\]
and
\[
mc(q_j(p))=c.
\]

Derive each firm's equilibrium price, quantity, and profit.



\subsection{Monopolistic Competition}

There is a relatively large group of firms selling differentiated products that buyers perceive as close, but not perfect, substitutes.

As a result, each firm has limited monopoly power. This makes monopolistic competition similar to the heterogeneous-goods Bertrand competition model.

A key difference, however, is that market entry is allowed.

When a new firm enters, it introduces a previously non-existent variety of the product.

Another difference is that firms produce very close substitutes.

In the short run, where \(J\) is fixed, we still have the equilibrium condition
\[
\frac{\partial q_j(p)}{\partial p_j}
\left[
mr(q_j(p))-mc(q_j(p))
\right]
=
0,
\]
where
\[
p_j<\bar p_j.
\]

Firms may have positive, negative, or zero short-run profit.

In the long run, firms exit the industry if their profits are negative. Positive long-run profits for any single firm induce the entry of new firms producing close substitutes.

As a result, in long-run equilibrium every active firm must have zero profit.

There are two long-run equilibrium conditions:
\[
\frac{\partial q_j(p^*)}{\partial p_j}
\left[
mr(q_j(p^*))-mc(q_j(p^*))
\right]
=
0,
\]
and
\[
\pi_j(q_j(p^*))=0.
\]

The zero-profit condition requires
\[
p_j^*
=
ac(q_j(p^*)).
\]

In the usual case where average cost \(ac(q)\) declines over the relevant output range,
\[
mc(q_j(p^*))<ac(q_j(p^*)).
\]

Thus the markup formula becomes
\[
\frac{
p_j^*-mc(q_j(p^*))
}{
p_j^*
}
=
\frac{
ac(q_j(p^*))-mc(q_j(p^*))
}{
p_j^*
}
=
\frac{1}{|\varepsilon_j(q(p^*))|}.
\]

This markup is still positive.

However, it cannot be large because
\[
|\varepsilon_j(q(p^*))|
\]
is large when products are close substitutes.

Because long-run production does not occur at the point where the firm's average cost
\[
ac(q_j(p))
\]
is minimized, it appears that there are too many firms in the industry.



in the industry, with each producing too little.

However, since products are not perfect substitutes, a greater number of firms increases variety for consumers.

As a result, it is unclear whether the long-run number of firms exceeds or falls short of the socially optimal level.

\subsection{Endogenous Number of Firms: Cournot Competition}

What will happen if entry and exit are free in an oligopoly market?

We consider the Cournot competition case with homogeneous products.

When \(J\) firms have entered the market, per-firm profit is
\[
\pi(J)
=
p(Jq(J))q(J)-c(q(J))-F,
\]
where \(F\) is a fixed setup cost.

The free-entry equilibrium number of firms, \(J^e\), satisfies the zero-profit condition:
\[
\pi(J^e)=0.
\]

The question is whether \(J^e\) is socially optimal.

The socially optimal number of firms, \(J^*\), solves, ignoring the integer constraint,
\[
\max_J W(J)
=
\int_0^{Jq(J)}p(x)\,dx
-
Jc(q(J))
-
JF.
\]

The first-order condition is
\[
W'(J)=0.
\]

Differentiating \(W(J)\) gives
\[
\begin{aligned}
W'(J)
&=
p(Jq(J))
\left[
q(J)+Jq'(J)
\right] \\
&\quad
-
c(q(J))
-
Jmc(q(J))q'(J)
-
F.
\end{aligned}
\]

Rearranging,
\[
W'(J)
=
\pi(J)
+
J
\left[
p(Jq(J))-mc(q(J))
\right]
q'(J).
\]

Thus, the first-order condition is
\[
\pi(J)
+
J
\left[
p(Jq(J))-mc(q(J))
\right]
q'(J)
=
0.
\]

A new entrant contributes directly to social welfare through its profit,
\[
\pi(J).
\]

The term
\[
J
\left[
p(Jq(J))-mc(q(J))
\right]
q'(J)
\]
is the additional indirect effect of altering the behavior of existing firms.

Since firms have monopoly power,
\[
p(Jq(J))-mc(q(J))>0.
\]

Because of the new entrant, existing firms may contract their output level, and thus
\[
q'(J)<0.
\]

Therefore, the indirect effect is negative.

This implies that at the socially optimal number of firms,
\[
W'(J^*)=0
\]
and
\[
\pi(J^*)>0.
\]

Because
\[
\pi(J^e)=0,
\]
and if \(\pi(J)\) decreases in \(J\), we conclude that
\[
J^e>J^*.
\]

Thus, symmetric Cournot competition with free entry results in excessive entry.

\paragraph{Exercise 5.}
Derive the conditions under which \(q(J)\) and \(\pi(J)\) decrease with respect to \(J\), respectively.

There are various models for estimating market power and performance.

Examples include Ackerberg et al.\ (2007) and Hortaçsu and Joo (2023).

Conventional methods, such as estimating demand systems or supply functions, have limitations because they assume that prices are exogenously given.

However, in a noncompetitive market, price must be treated as an endogenous variable, requiring econometric solutions to address endogeneity, for example, instrumental variables.

Furthermore, conventional approaches do not easily allow certain counterfactual analyses, such as how the market would respond to the introduction of a new product.

Another challenge is accounting for individual heterogeneity in product demand when estimating demand systems with market-level data.








